\clearpage
\begin{flushright}
  \fechaclase{28 de agosto de 2026}
\end{flushright}

\paragraph{Corolario}
La matriz
\[
  S=
  \begin{bmatrix}
    S_{11} & S_{12}\\
    S_{21} & S_{22}
  \end{bmatrix}
\]
es definida negativa si y solo si
\[
  S_{11}<0
\]
\[
  S_{22}<0
\]
\[
  S_{11}-S_{12}S_{22}^{-1}S_{12}^{T}<0
\]
\[
  S_{22}-S_{12}^{T}S_{11}^{-1}S_{12}<0
\]

\paragraph{\tituloejemplo{Ejemplo}}
Transforme la desigualdad de Lyapunov en una LMI
\[
  \underbrace{XA+A^{T}X}_{S_{11}}<0,
  \qquad A,X\in\mathbb{R}^{n\times n}
\]
\[
  S_{11}-S_{12}S_{22}^{-1}S_{12}^{T}<0
\]
\[
  \underbrace{XA+A^{T}X}_{S_{11}}
  -\underbrace{0_{n\times n}}_{S_{12}}
  \underbrace{I_{n\times n}^{-1}}_{S_{22}^{-1}}
  \underbrace{0_{n\times n}^{T}}_{S_{12}^{T}}<0
\]
\[
  \Rightarrow
  \begin{bmatrix}
    XA+A^{T}X & 0_{n\times n}\\
    0_{n\times n} & -I_{n\times n}
  \end{bmatrix}<0
\]

\[
  XA+A^{T}X+0I^{-1}0<0
\]
\[
  \underbrace{-(XA+A^{T}X)}_{S_{11}}
  -\underbrace{0}_{S_{12}}
  \underbrace{I^{-1}}_{S_{22}^{-1}}
  \underbrace{0}_{S_{12}^{T}}>0
\]
\[
  \Rightarrow
  \begin{bmatrix}
    -(XA+A^{T}X) & 0_{n\times n}\\
    0_{n\times n} & -I_{n\times n}
  \end{bmatrix}>0
\]

\paragraph{\tituloejemplo{Ejemplo}}
Convertir la ecuación de Riccati en una LMI
\[
  XA+A^{T}X^{T}+XBR^{-1}B^{T}X^{T}+Q<0
\]
con la solución $X=X^{T}$, $X\in\mathbb{R}^{n\times n}$
\[
  \underbrace{XA+A^{T}X^{T}}_{S_{11}}
  +\underbrace{XB}_{S_{12}}
  \underbrace{R^{-1}}_{S_{22}^{-1}}
  \underbrace{B^{T}X^{T}}_{S_{12}^{T}}+{Q}<0
\]
\[
  S_{11}-S_{12}S_{22}^{-1}S_{12}^{T}<0
\]
\[
  XA+A^{T}X^{T}
  -[XB(-R^{-1})B^{T}X^{T}]<0
\]
\[
  \Rightarrow
  \begin{bmatrix}
    XA+A^{T}X^{T} & XB\\
    B^{T}X^{T} & -R
  \end{bmatrix}<0
\]

\paragraph{\tituloejemplo{Ejemplo}}
Convertir la ecuación algebraica en una LMI
\[
  \underbrace{AX+X^{T}A^{T}-(BF+F^{T}B^{T})}_{S_{11}}
  +\underbrace{X}_{S_{12}}
  \underbrace{Q}_{S_{22}^{-1} | S_{22}}
  \underbrace{X^{T}}_{S_{12}^{T}}
  +FRF^{T}<0
\]
\[
  S_{11}-S_{12}S_{22}^{-1}S_{12}^{T}<0
\]
\[
  \underbrace{
    \begin{bmatrix}
      AX+X^{T}A^{T}-(BF+F^{T}B^{T}) & X\\
      X^{T} & -Q^{-1}
  \end{bmatrix}}_{S_{11}}
  +\underbrace{F}_{S_{12}}
  \underbrace{R}_{S_{22}^{-1}}
  \underbrace{F^{T}}_{S_{12}^{T}}<0
\]
\[
  \begin{bmatrix}
    AX+X^{T}A^{T}-(BF+F^{T}B^{T}) & X & F\\
    X^{T} & -Q^{-1} & 0\\
    F^{T} & 0 & -R^{-1}
  \end{bmatrix}<0
\]

\paragraph{\tituloejemplo{Ejemplo}}
Encuentre una solución $P=P^{T}\in\mathbb{R}^{2\times2}$ para la desigualdad
\[
  A^{T}P+PA<0,
  \qquad P>0,
  \qquad
  A=
  \begin{bmatrix}
    0 & 1\\
    -1 & -2
  \end{bmatrix}
\]
con el complemento de Schur
\[
  \begin{bmatrix}
    A^{T}P+PA & 0\\
    0 & -I
  \end{bmatrix}<0
\]

Con Sedumi/YALMIP
\begin{MatlabCode}
  clear; clc;

  P = sdpvar(2,2);
  A = [0 1; -1 -2];
  LMI1 = A'*P + P*A;
  LMI2 = P;

  cons = [LMI1 <= 0 , LMI1 >=0];
  solvesdp(cons)
  Pe = double(P)
  eig(Pe)

  eig(A'*Pe + Pe*A')
\end{MatlabCode}

\[
  V=X^{T}PX>0,
  \qquad \dot{X}=AX
\]
\begin{align*}
  \dot{V}
  &=X^{T}P\dot{X}+\dot{X}^{T}PX\\
  &=X^{T}PAX+X^{T}A^{T}PX\\
  &=X^{T}[PA+A^{T}P]X<0
\end{align*}

Ahora probando con
\[
  A=
  \begin{bmatrix}
    0 & 1\\
    -1 & 2
  \end{bmatrix}
\]
